\documentclass{article}
\usepackage{amsmath,hyperref}
\begin{document}
\title{Late evidence of one erased brake slot can require quadratically more formation clearance}
\maketitle

\paragraph{Physical model and evidence.}
Two robots translate in one dimension with initial equal speed
\(v\in\mathbb{N},v\ge1\), and must both stop, without
reversing, by cut \(n+1\), where \(n=2v\).
One unit of braking is available per unit slot, granted
exclusively to one robot. Action zero is allowed.
The initial gap lies strictly between a clearance and a
separation threshold; the change in gap \(q\) must stay in
\([-G,G]\) throughout. A separation threshold is a modeled
geometric restriction, not a radio guarantee.
At most one commanded slot can be omitted completely.
At slot \(e+d+1\), before its action, a source-qualified
execution receipt names the omitted slot \(e\); \(d=0\)
means notice before the very next action. Before receipt
there is no usable observation of the effect. The no-loss
case is admitted. The receipt has declared perfect naming
and delay in this model, not native hardware provenance.
A full grant to a stopped robot either clips to no effect
or is forbidden, as specified separately. The former is
a physical actuator behavior, not new information for the
controller. Scale length results by \(a\Delta^2\) for
single-robot brake acceleration \(a\) and slot duration
\(\Delta\), with initial dimensional speed \(va\Delta\).

With relative speed \(w\) and \(r=2q\), an effective
grant updates \(w'=w+s\), \(r'=r+2w+s\),
where \(s\in\{-1,1\}\); an omission, idle, or clipped
grant has \(s=0\). These are actual effects in the model,
not transmitted commands. Within a slot the relative
speed is affine with integer endpoints separated by at
most one, so the extrema of position change are attained
at endpoints. This gives exact continuous gap checks
using integer arithmetic.

\paragraph{Exact deadline and delay boundary.}
The no-loss path must use each sign exactly \(v\) times
in its first \(n\) slots: a fault can erase the last slot.
At the cut before slot \(n\), a loss in any of the last
\(\min(d,n)\) planned useful slots is still hidden. If one
of them was omitted, the final grant must name precisely
the robot assigned to that slot. All these hidden slots
of the no-loss schedule must therefore have the same sign.
If \(d>v\) this is impossible even with clipping and no
finite gap bound. If clipping is forbidden, for every
\(d\ge1\) the hidden loss of the previous grant requires
a final action the no-loss world forbids, also independent
of the gap bound. This is not a lack of physical capacity:
all worlds are individually solvable with timely evidence.

For \(1\le d\le v\), a nominal balanced prefix of
\(v-d\) pairs followed by \(d\) positive and \(d\)
negative grants supplies an unconditional upper
construction. Continue the nominal schedule until a receipt
arrives; insert the named lost grant then resume the nominal
sequence, shifted one slot. If the loss remains hidden
at the final cut, grant to the common sign of the last
\(d\) slots. It repairs a lost grant or clips to no effect.
Thus all velocities stop on time. This uses only the scheduled
grants and the delivered name of the lost slot, never a hidden
actual velocity.

\paragraph{Quadratic cost of delayed evidence.}
Choose sign so that the last \(d\) nominal grants reduce
relative speed from \(d\) to zero. On their first slot an
unobserved omission preserves speed \(d\) and increases
\(q\) by \(d\). The later \(d\) useful grants add \(d^2/2\)
before stop, independently of the rest of the policy.
The required band therefore has
\(G\ge d^2/4+d/2\).
The invariant \(r+w^2\) under the useful tail grants
gives this area exactly. There is also \(G\ge2\)
for any winning policy under one omission: after the
first signed grant choose the second slot to be omitted,
forcing \(q=1.5,w=1\) up to sign at the next cut; even
the fastest reduction of \(w\) then forces \(q\ge2\).

In the constructed balanced-prefix policy the nominal
gap is bounded by \(d^2+1\) and nominal relative
speed by \(d\). During the unobserved interval, omitting a
grant shifts the nominal gap by at most \(d+\tfrac12\).
Insert the lost grant at the receipt and let \(w_t\)
be the nominal speed difference where the insertion occurs.
On rejoining the nominal schedule one slot late, the
constant displacement offset is
\(w_t-\sigma(d+1)\), with \(\sigma\) the omitted sign.
Its magnitude is at most \(2d+1\). The final blind
compensation of a still hidden late loss is obeyed with
a no-larger offset. Thus \(G\le d^2+2d+2\), independent
of \(v\). This proves a \(\Theta(d^2)\) slack cost for
\(1\le d\le v\) at the one-spare-slot deadline, not just
a sufficiency inequality.

\paragraph{Further exact slices.}
For \(d=0\) repeat the named lost grant immediately
and then continue alternating ordered balanced pairs.
Nominal gap oscillates between 0 and 1; one erased slot
can shift that by at most 1. Stop even without clipping
by idling if the original motion already completed.
For \(d=1\) with clipping, use the same nominal pattern
\(+, -,-,+\) repeated, possibly ending with one pair.
At first notice of a nonfinal omission, the relative
state just after the next planned action, up to its
corresponding residue of \(e\) mod \(4\), is
\[
(q,w)=(-.5,-1),\ (2,0),\ (1.5,1),\ (-1,0).
\]
For even residues immediately compensate the omitted sign.
For odd residues compensate it and then perform the
opposite grant. These corrections leave equal remaining
counts and zero relative velocity, at gaps \(-1,1,2,0\),
respectively. Use subsequent balanced pairs toward zero.
If the last nominal grant is omitted but still unknown,
its sign is the final default; if no omission occurred
that action clips. This achieves \(G=2\) for all \(v\),
matching the lower bound.

At \(d=v\) the nominal order is necessarily two
opposite homogeneous blocks of size \(v\).
An omission at the start of the second block yields
gap \(v^2+v\) up to sign. The blocks plus final correction
attain exactly that symmetric bound. For \(v\ge3,d=v-1\),
take the last \(v-1\) signs to be negative. The only
negative sign in the prefix is at index \(j\in[0,v]\).
Nominal terminal gap is \(v^2-2v+2j\); a loss at
the first slot of the tail adds \(v-1\). Hence
\(G\ge v^2-v-1\). The choice \(j=0\) attains it.
If the first negative grant is omitted, its receipt
arrives after \(v-1\) effective positive grants;
then use \(v\) negatives followed by a positive, whose
peak gap is \((v-1)^2\).
A loss in the positive block can decrease the nominal
gap but never below \(-2\); by the end its gap is
\(v^2-4v+j'\) for the lost index \(j'\in[1,v]\).
A loss in the final negative block adds at most \(v-1\).
These cases cover the no-loss history as well.
Intermediate delay values are not assigned an exact
closed form by these slices.

\paragraph{What changing the deadline buys.}
With no feedback at all and clipping allowed, two spare slots
rather than one are necessary and sufficient for deadline
feasibility. Each sign needs at least \(v+1\) scheduled
opportunities, since either sign can lose one. Balanced
nominal pairs plus one extra of each sign provide that budget.
For exactly \(2v+2\) slots with no feedback, the final
gap shift upon erasing the first grant of a sign is the
signed span between its first and last scheduled slots:
one speed unit remains until its extra grant, including the
two half-slot ramps. The two signs shift in opposite
directions, each by at least \(v\), requiring \(G\ge v\).
The balanced construction has \(G\le2v+2\).
Thus a looser deadline repairs feasibility, not the
geometry cost of retaining no execution evidence.

\paragraph{Comparison and residuals.}
Full causal finite enumeration agrees on the checked
small instances. Separately, the analytical policies are
replayed with explicit indistinguishable observation
prefixes and actual effects. Where two half grants can
physically actuate together, equal braking over each
good slot stops both at the same deadline with no gap
change and no receipt, given clipping. That is a stronger
actuator, not an implementable comparison for the
declared exclusive one. Breeden and Panagou's
\href{https://arxiv.org/html/2210.01354v2}{barrier-composition model}
assumes a convex control set in Section II-B; the
exclusive unit-slot actuation here is outside that
premise. The result extends the no-omission shared-brake
slice in \#611 but neither provides a receipt sensor
nor claims a robot control or independent mission
result. It supplies a quantified R6 separation among
gap feasibility, timed stopping, available observations,
and actual authority over the shared resource.
\end{document}
